% Chapter 7 — Conclusion and Future Work

\chapter{Conclusion and Future Work}
\label{chap:conclusion}

\section{Conclusion}

This thesis implemented and compared a range of regret-minimizing algorithms
within a common experimental framework. The evaluation covered external,
internal, and swap regret under both full-information and bandit feedback, as
well as repeated-game behavior in RPS and RPSLS.

The experiments show that the choice of regret notion is important for an
empirical comparison. Differences that are small under external regret can
become much clearer once action-dependent deviations are considered. At the
same time, the results depend on the reward process, action-space size,
horizon, and, in repeated games, the opponent. The experiments therefore do
not suggest a single ordering of the algorithms that applies across all tested
settings.

The repeated-game experiments also connect these regret differences to
equilibrium behavior. External regret and stronger regret notions can lead to
different CE and CCE behavior, while cross-play shows that the observed regret
of a learner can depend strongly on the learning dynamics of its opponent.
Together with the finite-horizon scaling experiments, this provides an
empirical view of distinctions that are less visible from asymptotic regret
bounds alone.

The theoretical comparison in Chapter~\ref{chap:discussion} also shows where
the common action-regret evaluator is informative and where it is not.
Relations between action regret, distribution regret, and pseudo-regret allow
some empirical conclusions to be transferred between metrics, while
sub-root mixed-strategy rates and stronger probabilistic guarantees cannot be
resolved directly with the current evaluation.

% ============================================================
\section{Future Work}

A natural next step is to add regret evaluators that more closely match the
quantities used in the theoretical analyses. In particular, distribution and
fully mixed regret would allow a more direct comparison with the guarantees
for Hedge-based and optimistic methods. For high-probability results, replicate
quantiles or empirical tail frequencies would complement the mean regret used
here.

The repeated-game experiments could also be extended beyond the selected RPS
and RPSLS profiles. A larger cross-play matrix and a broader set of games would
make it possible to examine whether the opponent effects observed here persist
in other settings. Equilibrium evaluation could also report the largest remaining
profitable deviation, in addition to the geometric distance to the CE and CCE sets.

Finally, longer horizons, more replicates, and larger action spaces would help
separate persistent scaling behavior from finite-horizon effects. These
extensions would use the same experimental framework while providing a more
complete comparison between theoretical regret guarantees and observed
learning behavior.
